% STATUS: written, 20 October 2026.
%
% Most of this is transcribed from docs/mdp-specification.md and
% docs/adr/ADR-0002-flat-declarative-action-space.md. If either of those
% changes, change this with it.

\section{System Model and Methodology}\label{sec:model}

\subsection{Placement decisions and their cost}

A cluster has $R$ regions, each holding one database node. The keyspace is
divided into $C$ atomic cells, and a \emph{shard} is a contiguous run of cells.
Each shard is stored as several \emph{replicas}, which hold the same data. One of them is the \emph{primary}, which CockroachDB calls the \emph{leaseholder}. The primary serves writes and, because this work uses strong reads throughout, all reads. The number of replicas per shard, called the replication factor, is fixed at three throughout this work, and the agent does not control it.

Three decisions are open at any moment: which region holds a
given shard's primary, which regions hold its replicas, and where the boundary
between two shards falls. The first two are allocation and the third is
fragmentation. We treat them as one problem because they interact:
splitting a shard is only worth doing if the two halves are then placed
differently.

Latency is computed analytically from a fixed matrix of round-trip times
between regions. A read costs the round trip to the shard's primary, because on the real cluster a strong read (one that must see the latest committed write) is served there, and the simulator must price it the same way. A write costs
the round trip to the primary plus the time for a majority of replicas to
acknowledge it, since a write is not durable until enough copies agree. A
transaction touching several shards costs its member operations plus, if those
shards' primaries sit in more than one region, an additional charge for the
two-phase commit (2PC) protocol that keeps it atomic. This work depends on
that last term. It is what rewards placing the primaries
of two shards that are used together in the same region, and so it is what
links the fragmentation decision to the allocation decision.

\subsection{The workload}

The workload generator produces a stream of requests whose geographic origin
moves. Each of the $C$ cells sits at an angular position on a circle, as do the
regions, and traffic for a cell comes mostly from the regions near its angle.
The whole angular field rotates over a daily period, which is what makes a
placement chosen at one time wrong at a later one. Cell popularity follows a
Zipf distribution, so a minority of cells carry most of the traffic, and
localized spikes, each of which multiplies one region's traffic for a short
burst, arrive at random times (a Poisson process). The read and write mix and
the Zipf skew follow the style of YCSB, the Yahoo!\ Cloud Serving
Benchmark~\cite{cooper2010ycsb}, though our values (read fraction 0.8, Zipf
exponent 0.9) match none of its core workloads; geographic origin is the
dimension YCSB does not model. Adjacent cells get correlated but not identical
angular positions from a random walk, which is what makes splitting
worthwhile: a wide shard spans divergent angles and cannot be placed well
anywhere, whilst its two halves can. A single parameter, \emph{volatility},
scales the rotation speed, the spike rate and the spike size together, from
slow drift with rare, small spikes to fast drift with frequent, large ones; it
is the main setting the evaluation varies.

\subsection{The decision problem}

We state the problem as a Markov decision process: once per decision window
the agent observes the state, takes one action, and receives a reward.

\paragraph{Observation} The state is encoded as a fixed-length vector
describing, for every shard, its read, write and transaction demand by
region of origin, its size, and the regions holding its primary and replicas;
for every region, its share of requests and its storage occupancy;
the recent history of data movement; the matrix of inter-region round-trip
times; and a small set of scalars summarizing current cost.

\paragraph{Actions} The action space is flat and declarative, with 261 actions
at $R=6$ regions and a cap of 20 shards. There are four kinds of action:

\begin{itemize}
\item \textsc{move\_lease}$(f, r)$ makes region $r$ the primary for shard $f$.
It moves no data and is only legal where $r$ already holds a replica.
\item \textsc{move\_replica\_to}$(f, r)$ places a replica of shard $f$ in
region $r$, evicting the non-primary replica furthest from $r$ so that the
replication factor stays constant. It costs one shard's worth of movement.
\item \textsc{split}$(f)$ divides shard $f$ at its midpoint.
\item \textsc{noop} does nothing.
\end{itemize}

Every action names a desired end state, not an operation, so every action is
defined in every state: one that cannot be applied costs a small penalty instead
of raising an error, and we also mask such actions out for the agent
(Section~\ref{sec:agent}). We chose this encoding by measuring the
alternatives: the destination-only form of \textsc{move\_replica\_to} gives 261
actions against 861 for the naive source-and-destination form, more of its
actions apply in a typical state, and a breadth-first search over the
reachability graph confirms that no placement becomes unreachable.

\paragraph{Reward} The reward is the negative weighted sum of three normalized
terms: mean latency divided by a scale constant, data moved as a fraction of
the total dataset, and load imbalance. Movement is a fraction of the dataset,
not bytes, because the simulator holds 60\,GB whereas the real cluster holds
about 1\,GB; a reward in bytes would score the same physical move sixty times
differently in the two environments.

\subsection{Baselines}

We compare against six reference points: five policies, one of which (CockroachDB) runs only on the cluster, and an oracle.

\begin{itemize}
\item \textbf{Static hash} holds, for the whole episode, the placement a
consistent-hash partitioner~\cite{karger1997consistent} produces, so it never
moves anything. It is the floor any policy must clear.
\item \textbf{Random} selects an action uniformly from the whole space. It is a control: it
tests whether any gain comes from choosing \emph{where} to move data rather than
from moving data at all.
\item \textbf{Follow traffic} finds, each window, the single shard whose
primary placement costs the most relative to the best available region, and
moves it towards that region. It is the idea behind CockroachDB's built-in
follow-the-workload lease rebalancing~\cite{crdb2025ftw}, applied every window
and to replicas as well as leases. It looks only at the last window: it weighs
no migration cost and anticipates no drift.
\item \textbf{Lion-style} is a co-location planner after
Lion~\cite{lion2024}. Every 10 windows it groups shards that are accessed
together, chooses a region for each group with a cost model that prefers a
lease move (cheap, since it only changes which replica is primary) over a replica move (which copies a shard), and then
rebalances load. It is the closest published system we could adapt to this
action space, and unlike the heuristic it weighs movement cost. Like Lion, its
cost model has no distance term, and we omit Lion's workload predictor, so it
plans from the last window; it does nothing in about 89\% of windows.
\item \textbf{CockroachDB} hands placement back to the database, on the
cluster only. From the common starting placement the harness re-enables its
load-based lease and replica rebalancing~\cite{crdb2025settings}, follow-the-workload included~\cite{crdb2025ftw}, and then chooses nothing. Splits and merges stay
off, as for every policy. The simulator has no model of it.
\item \textbf{Per-window oracle (fixed shards)} is not a policy. Each window,
it searches every assignment of primaries to regions ($6^6$ for the six
starting shards), with the other replicas in the regions nearest each primary, and
keeps the one with the lowest latency under the full cost model, commit
coordination included. Movement is free and the window is seen in advance.
Under strong reads, commit coordination is the only cost that links one shard's
placement to another's, so the search is exact. Its result is a lower bound on the latency of any placement of the
starting shards; only a policy that splits can score below it.
\end{itemize}

\subsection{The agent}\label{sec:agent}

The agent is a Deep Q-Network. A Q-network is a function that, given the
current state, estimates the long-run value of each available action; the
policy is then to take the action with the highest estimate. The network here
has two hidden layers of 256 units and is trained with the standard
Stable-Baselines3 implementation~\cite{raffin2021sb3} against a Gymnasium
environment~\cite{towers2025gymnasium}, following Mnih et
al.~\cite{mnih2015dqn}. The learning rate is $10^{-4}$, the replay buffer holds
$10^5$ transitions, the batch size is 128, and the discount is 0.99. The target network
(a lagged copy of the network, used to compute training targets) is refreshed every 2000 steps,
and the exploration rate $\epsilon$ (the chance of taking a random action) falls from 1.0 to 0.05 over the
first 30\% of training. Every $10^4$ steps the policy is scored on three
validation seeds (1000 to 1002), and the checkpoint with the lowest mean latency is kept.

We add two things to the plain Deep Q-Network. First, we mask inapplicable
actions~\cite{huang2022masking,zahavy2018elimination}, both when the agent
selects an action and when it computes its training target (the highest value
it predicts for the state the return ends in). In earlier runs under a
different protocol, across several volatilities and migration weights, 79.9 to
89.8\% of an unmasked agent's actions were inapplicable, no better than a
uniform random policy's 80.4\%. Second, returns are summed over three steps
rather than one, because a move is charged in full when it is made whereas the
latency it saves arrives in later windows. Under the protocol of
Section~\ref{sec:metrics}, three learner seeds trained both ways differ by
0.3\,ms on average (95\% interval $[-3.0, 2.3]$\,ms over workload seeds); an
earlier comparison at replication factor 2 found a 5.4\,ms improvement over
four learner seeds that was not significant (paired 95\% $t$-interval
$[-14.6, 3.8]$\,ms). We kept three-step returns because they cost nothing. We
use none of the other extensions that Rainbow~\cite{hessel2018rainbow} combines
(double Q-learning, a dueling head, prioritized replay, distributional value
learning and noisy networks), so that any later
improvement can be traced to the extension that produced it.

Training cycles through 64 workload seeds, not one, because an agent trained on
a single workload learns that workload instead of the pattern behind it.
Evaluation is greedy (no exploration) and uses seeds disjoint from the training
seeds.

\subsection{Metrics}\label{sec:metrics}

Results report mean latency in milliseconds, averaged over the windows of an
episode (one run of 200 decision windows); the 95th-percentile latency (p95)
where the tail matters; data moved, in gigabytes per episode; and load
imbalance, how unevenly requests are spread across regions. Simulator results
are means over eight workload seeds. Where a comparison turns on a difference,
we give a 95\% percentile bootstrap~\cite{efron1979bootstrap} interval over
those seeds, resampling the per-seed differences with replacement, because
eight is a small sample and nothing guarantees the per-seed means are normally
distributed.

\subsection{Run records and the real cluster}

One script writes a self-contained directory per simulator run, holding its
configuration, seeds, code commit, package versions, machine, raw per-window
record, per-seed summary and the command that reproduces it; the simulator is deterministic, so a run reproduces
exactly. Training and cluster runs execute unattended on GitHub Actions, each
recording the hardware it received. Agent and cluster figures are copied from
those runs' summaries and archived artifacts, and each recorded table names
the run it came from. \ifanonymous An anonymized copy of the code, configurations, run records and archived artifacts accompanies the submission.\else The code, configurations, run records and archived artifacts are in the project repository.\fi

The trained agent runs on the cluster unchanged. Training in simulation and
deploying on hardware is borrowed from robotics~\cite{tobin2017domain}, but here
the gap between the two is narrower and easier to check, because
the simulator and the cluster can be built to agree on the observation space,
the action space and the cost structure. We deploy six
CockroachDB~\cite{taft2020cockroachdb} nodes, one per simulated region, in
containers on a single host, with wide-area round-trip times injected between
them. We chose CockroachDB from more than twenty systems we surveyed because it
offers a \emph{synchronous} statement, one that returns only after the change
has taken effect, for each of the agent's three actions. Synchrony is required:
if a placement command returned before the placement changed, the reward
observed after an action would belong to a state the action did not produce,
and the decision problem would no longer be well defined.

\begin{figure}[t]
\centering
\begin{tikzpicture}[
  node distance=3mm,
  box/.style={draw, rounded corners=1pt, align=center, font=\scriptsize,
              minimum width=30mm, inner sep=1.6mm},
  title/.style={font=\scriptsize\bfseries},
  parity/.style={<->, >={Stealth[length=1.6mm]}, dashed, gray},
  note/.style={font=\tiny, gray, align=center}
]

% Left column: the simulator.
\node[title] (simtitle) {Simulator};
\node[box, below=2mm of simtitle] (simwork) {drifting workload\\256 cells, 6 regions};
\node[box, below=of simwork] (simcost) {analytic cost model\\reads, writes, 2PC};
\node[box, below=of simcost] (simobs) {observation encoder};

% Right column: the real cluster.
\node[title, right=26mm of simtitle] (clutitle) {CockroachDB cluster};
\node[box, below=2mm of clutitle] (cluwork) {load generator\\6 nodes, netem latency};
\node[box, below=of cluwork] (clucost) {measured latency\\from real transactions};
\node[box, below=of clucost] (cluobs) {observation encoder};

% The three parity constraints.
\draw[parity] (simwork) -- node[note, above] {action space} node[note, below] {261 actions} (cluwork);
\draw[parity] (simcost) -- node[note, above] {cost structure} (clucost);
\draw[parity] (simobs) -- node[note, above] {observation space} (cluobs);

% The agent, underneath, moving left to right.
\node[box, below=6mm of $(simobs)!0.5!(cluobs)$, minimum width=62mm]
  (agent) {Deep Q-Network\\trained on the left, evaluated unchanged on the right};
\draw[->, >={Stealth[length=1.6mm]}] (simobs) -- (agent);
\draw[->, >={Stealth[length=1.6mm]}] (cluobs) -- (agent);

\end{tikzpicture}
\caption{The two halves of the system and the three things that must agree
between them. The agent is trained against the simulator, which is fast and
cheap, and evaluated without further training on the cluster, which is slow and
real. Transfer is
only meaningful if the observation space, the action space and the cost
structure are the same on both sides.}
\label{fig:architecture}
\end{figure}
